% STATUS: DRAFT
\chapter{The crossed product construction}\label{ch:crossed}

\begin{physmotiv}
Modular flow is the hidden dynamics of a state, but on a type III algebra it is \emph{outer}: no operator of the algebra generates it (\cref{ch:modham}). The cure is to enlarge the algebra by the missing generator. Concretely, we adjoin a new degree of freedom (a ``clock'' $x$ with conjugate momentum $p$) and an operator $X=K+x$ whose adjoint action is the modular flow. In the enlarged algebra the flow is \emph{inner} by construction. This is the crossed product. It is a mathematical device in Part V; in gravity (Part VII) the clock is a physical observer and the enlargement is forced by the constraints.
\end{physmotiv}

\section{Abstract definition}

Let $\alpha:\R\to\Aut(\M)$ be a $\sigma$-weakly continuous one-parameter group of automorphisms of $\M\subset\BH$. On $L^2(\R,\HH)=\HH\otimes L^2(\R)$ define
\begin{equation}
(\pi_\alpha(a)\xi)(s)=\alpha_{-s}(a)\,\xi(s),\qquad(\lambda(t)\xi)(s)=\xi(s-t).
\end{equation}
Then $\pi_\alpha$ is a faithful representation of $\M$, $\lambda$ is a unitary representation of $\R$, and they satisfy the \emph{covariance relation}
\begin{equation}
\lambda(t)\,\pi_\alpha(a)\,\lambda(t)\ad=\pi_\alpha\big(\alpha_t(a)\big).
\end{equation}
(Check: both sides act on $\xi(s)$ by $\alpha_{t-s}(a)$.)

\begin{definition}
The \emph{crossed product} is the von Neumann algebra $\M\rtimes_\alpha\R=\{\pi_\alpha(a),\lambda(t):a\in\M,t\in\R\}''$ on $\HH\otimes L^2(\R)$.
\end{definition}

In words: $\M\rtimes_\alpha\R$ is generated by $\M$ and a unitary group $\lambda(t)$ that implements $\alpha_t$. It is the smallest algebra containing $\M$ in which $\alpha$ becomes inner. Elements have the form $\int\dd t\,a(t)\lambda(t)$ with $a(t)\in\M$, composed by the twisted convolution $(a\star b)(t)=\int\dd s\,a(s)\alpha_s(b(t-s))$.

\section{The physicist's form}

Take now $\alpha=\sigma^\Omega$, the modular flow of a cyclic separating vector $\Omega$. Put $K=-\log\Delta_\Omega$, so $\sigma_t(a)=\ee^{-\ii tK}a\,\ee^{\ii tK}$ and $K\Omega=0$. Let $x$ be the position operator on $L^2(\R)$ and $p=-\ii\,\dd/\dd x$.

\begin{proposition}
The crossed product is unitarily equivalent to
\begin{equation}
\M\rtimes_\sigma\R\;\cong\;\big\{\,a\otimes\id\ (a\in\M),\ \ X\equiv K+x\,\big\}''\ \subset\ \Bnd{\HH\otimes L^2(\R)}.
\end{equation}
Here $\ee^{-\ii tX}$ plays the role of $\lambda(t)$, and $X$ is affiliated with the algebra.
\end{proposition}
\emph{Check.} Since $x$ commutes with $a\otimes\id$, $\ee^{-\ii tX}(a\otimes\id)\ee^{\ii tX}=\ee^{-\ii tK}a\ee^{\ii tK}\otimes\id=\sigma_t(a)\otimes\id$. So $\ee^{-\ii tX}$ implements the modular flow, as required. (The identification with the $\pi_\alpha,\lambda$ picture is a unitary change of frame.)

\begin{keyidea}
In $\M\rtimes_\sigma\R$ the modular flow is \emph{inner}: it is $\Ad\,\ee^{-\ii tX}$ with $X=K+x$ in the algebra. The price is the new variable $x$. Note $K$ alone is \emph{not} in the algebra; only the combination $K+x$ is.
\end{keyidea}

\subsection*{Second frame (the one used in gravity)}
Conjugating by $\ee^{-\ii pK}$ and using $\ee^{-\ii pK}x\,\ee^{\ii pK}=x-K$ (since $\ee^{\ii cp}x\ee^{-\ii cp}=x+c$), we get $\ee^{-\ii pK}(K+x)\ee^{\ii pK}=x$ and
\begin{equation}
\M\rtimes_\sigma\R\;\cong\;\big\{\,\ee^{-\ii pK}a\,\ee^{\ii pK}\ (a\in\M),\ \ x\,\big\}''=\{\sigma_p(a),\,x\}'',\label{eq:frame2}
\end{equation}
where $\sigma_p(a)$ means the modular flow with the \emph{operator} $p$ as time. This is the form in which the algebra arises as the set of gauge-invariant observables in gravity (\cref{ch:add_observer,ch:clpw}): $x$ becomes (minus $\beta$ times) the observer's energy and $\sigma_p(a)$ is a field operator dressed to the observer. The two frames contain the same information; frame~1 is better for the mathematics, frame~2 for the physics.

\section{The commutant}

In frame 1 the commutant is again a crossed product:
\begin{equation}
(\M\rtimes_\sigma\R)'=\big\{\,\ee^{\ii pK}(a'\otimes\id)\ee^{-\ii pK}\ (a'\in\M'),\ \ x\,\big\}''.
\end{equation}
\emph{Check.} (i) $\ee^{\ii pK}a'\ee^{-\ii pK}=\Delta^{-\ii p}a'\Delta^{\ii p}\in\M'$ (the modular flow preserves $\M'$ too), so it commutes with $a\otimes\id$. (ii) $x$ commutes with $a\otimes\id$ and with $X=K+x$. (iii) Using $\ee^{-\ii pK}X\,\ee^{\ii pK}=x$, we get $[X,\ee^{\ii pK}a'\ee^{-\ii pK}]=\ee^{\ii pK}[x,a']\ee^{-\ii pK}=0$. That these generate the \emph{whole} commutant is part of Takesaki's theorem. In particular \emph{the commutant of the crossed product of the right wedge is (the crossed product of) the algebra of the left wedge}, again with the same variable $x$.

\section{Two sanity checks}

\subsection{Inner flows: the crossed product is trivial}
Let $\M=M_n$ on $\HH=\C^n\otimes\overline{\C^n}$ and $\Omega=\rho^{1/2}$, so $K=K_\M-K_{\M'}$ with $K_\M=-\log\rho\in\M$ (\cref{ch:tomita}). Then $y\equiv X-K_\M=x-K_{\M'}$ commutes with $\M$, with $K_\M$ and with $x$, and $X=K_\M+y$ with $K_\M\in\M$. Hence the algebra generated by $\M$ and $X$ is the algebra generated by $\M$ and $y$:
\begin{equation}
M_n\rtimes_{\sigma}\R\;\cong\;M_n\otimes\Linf(\R_y).
\end{equation}
Type I $\otimes$ abelian, with centre $\Linf(\R_y)$: the crossed product by an inner flow adds nothing but a classical variable. In particular it is of type I$_n$ over the centre and not a factor.

\subsection{The infinite spin chain}
For the Powers factor $R_\lambda$ with its periodic modular flow, $\sigma_{t_0}=\mathrm{id}$ for $t_0=2\pi/|\log\lambda|$. Then the crossed product has a nontrivial centre (functions of $\ee^{\ii t_0 X}$-type), reflecting that $\R/t_0\Z$-dual circle survives. For III$_1$ the flow has no period and the centre is trivial (next chapter).

\begin{whatwrong}
\begin{itemize}
\item \textbf{State independence.} The crossed product $\M\rtimes_{\sigma^\varphi}\R$ does not depend on which faithful normal state or weight $\varphi$ is used: for another $\psi$ the Connes cocycle $(D\psi:D\varphi)_t$ (\cref{ch:modham}) gives an isomorphism. Hence \emph{the algebra $\M\rtimes\R$ is canonical}, although its explicit realization uses a state.
\item \textbf{Not a tensor product.} $\M\rtimes_\sigma\R\not\cong\M\bar\otimes\Linf(\R)$ in general. The twisting by $\sigma$ is the whole point.
\item \textbf{Which time?} We cross by the \emph{full} modular flow. Crossing by a different one-parameter group (e.g.\ a flow that is not modular) gives different, possibly type III, algebras.
\end{itemize}
\end{whatwrong}

\section*{Exercises}
\begin{exercise}
Verify the covariance relation $\lambda(t)\pi_\alpha(a)\lambda(t)\ad=\pi_\alpha(\alpha_t(a))$ for $\pi_\alpha,\lambda$ as defined. Check that $\pi_\alpha$ is a $^*$-homomorphism.
\end{exercise}
\begin{exercise}
Verify $\ee^{\ii cp}x\,\ee^{-\ii cp}=x+c$ and the identity $\ee^{-\ii pK}(K+x)\ee^{\ii pK}=x$. Deduce \eqref{eq:frame2}.
\end{exercise}
\begin{exercise}
Verify that $x$ and $\ee^{\ii pK}a'\ee^{-\ii pK}$ commute with $a\otimes\id$ and with $K+x$.
\end{exercise}
\begin{exercise}
For $\M=M_2$ with $\rho=\mathrm{diag}(p,1-p)$, write out $K_\M$, $K_{\M'}$ and $y$ explicitly, and exhibit the isomorphism $\M\rtimes\R\cong M_2\otimes\Linf(\R_y)$. Which element is the centre?
\end{exercise}
\begin{exercise}
Let $\alpha_t=\Ad\,\ee^{-\ii tH}$ for a bounded self-adjoint $H\in\M$ (an inner flow), and let $X$ be the generator of $\lambda$ in frame 1 (so $\ee^{-\ii tX}$ implements $\alpha_t$). Show that $y=X-H$ commutes with $\M$ and with $X$, and conclude $\M\rtimes_\alpha\R\cong\M\bar\otimes\Linf(\R_y)$.
\end{exercise}
